\exercisesection{\S~7}

所考察的所有 g-模（习题14和15中的模除外）均假定为有限维的。

\begin{enumerate}
\def\labelenumi{\arabic{enumi})}
\tightlist
\item
  令 \(\omega \in \mathrm{P}_{++}\)；以 \(\mathrm{S}(\omega)\) 表示 \(\mathrm{E}(\omega)\) 的权集，换言之，它是包含 \(\omega\) 的最小 R-饱和子集（命题5）。若 \(\lambda \in \mathrm{S}(\omega)\)，则有 \(\lambda \equiv \omega\)（mod. Q）。反之，令 \(\lambda \in \mathrm{P}\) 满足 \(\lambda \equiv \omega\)（mod. Q)；证明下列性质等价：
\end{enumerate}

\begin{enumerate}
\def\labelenumi{(\roman{enumi})}
\item
  \(\lambda \in S(\omega);\)
\item
  \(\omega - w\lambda \in Q_{+}\) 对所有 \(w \in W\) 均成立；
\item
  \(\lambda\) 属于 W 的 \(\omega\) 在 \(\mathfrak{h}_{\mathbf{R}}^{*}\) 中的凸包。
\end{enumerate}

（为证明 (iii) \(\Longrightarrow\) (ii)，注意 \(\omega - w\omega\) 是 \(\mathrm{R}_{+}\) 中元素的系数 \(\geq 0\) 的线性组合；由凸性推出 \(\omega - w\lambda\) 也有同样的性质；由于 \(\omega - w\lambda\) 属于 Q，这意味着 \(\omega - w\lambda \in \mathrm{Q}_{+}\)。为证明 (ii) \(\Longrightarrow\) (i)，选取 \(w\) 使得 \(w\lambda \in \mathrm{P}_{++}\)，并应用命题3的推论2。(i) \(\Longrightarrow\) (iii) 的蕴含是直接的。）

\begin{enumerate}
\def\labelenumi{\arabic{enumi})}
\setcounter{enumi}{1}
\tightlist
\item
  令 \((\mathrm{R}_{i})_{i\in\mathrm{I}}\) 为 R 的不可约分支族，且 \(g=\prod_{i\in I}g_{i}\) 为 g 分解为单代数之积的相应分解。将 P 与 \(\mathrm{P}(\mathrm{R}_{i})\) 的乘积相识别，并给每个 \(\mathrm{P}(\mathrm{R}_{i})\) 赋予由基 \(B_{i}=B\cap R_{i}\) 定义的序关系。
\end{enumerate}

\begin{enumerate}
\def\labelenumi{\alph{enumi})}
\item
  令 \(\omega = (\omega_{i})_{i\in \mathrm{I}}\) 为 \(\mathrm{P}_{++} = \prod_{i\in \mathrm{I}}\mathrm{P}_{++}(\mathrm{R}_i)\) 中的一个元素。证明单 \(\mathfrak{g}\)-模 \(\operatorname {E}(\omega)\) 同构于单 \(\mathfrak{g}_i\)-模 \(\operatorname {E}(\omega_i)\) 的张量积。
\item
  令 \(\mathcal{M}\)（分别地，\(M_{i}\)）为 \(P_{++}\)（分别地，\(P_{++}(R_{i})\)）中具有命题6和7的 (i)、(ii)、(iii)、(iv) 的等价性质的元素集。证明 \(M = \prod_{i \in I} M_{i}\)，换言之，\(\omega \in M\) 当且仅当对于所有 \(i \in I\)，\(\omega_{i}\) 或为零，或为 \(R_{i}\) 的极小权。推出 M 是 P/Q 中元素在 P 里的代表元系。
\item
  令 E 为单 g-模，X 为其权集。证明 X 含有 M 中唯一的一个元素，且该元素的重数等于 X 中各元素重数的上界。
\end{enumerate}

\begin{enumerate}
\def\labelenumi{\arabic{enumi})}
\setcounter{enumi}{2}
\tightlist
\item
  \begin{enumerate}
  \def\labelenumii{\alph{enumii})}
  \tightlist
  \item
    令 E 为 g-模。证明下列条件等价：
  \end{enumerate}
\end{enumerate}

\begin{enumerate}
\def\labelenumi{(\roman{enumi})}
\item
  \(\mathfrak{g}\) 与 E 的半直积的秩严格大于 \(\mathfrak{g}\) 的秩。
\item
  0 是 E 的一个权。
\item
  \(\mathrm{E}\) 存在一个根权（即~属于 \(\mathrm{Q}\) 的权）。
\end{enumerate}

\begin{enumerate}
\def\labelenumi{\alph{enumi})}
\setcounter{enumi}{1}
\tightlist
\item
  假设 \(\mathrm{E}\) 是单的。证明 (i)、(ii)、(iii) 等价于
\end{enumerate}

\begin{enumerate}
\def\labelenumi{(\roman{enumi})}
\setcounter{enumi}{3}
\tightlist
\item
  E 的最高权是根权。
\end{enumerate}

若这些条件满足，则不存在非零不变交错双线性型。（使用命题12和§6的命题1 (ii)。）

\begin{enumerate}
\def\labelenumi{\arabic{enumi})}
\setcounter{enumi}{3}
\item
  令 \(k'\) 为 \(k\) 的一个扩张，且 \(\mathfrak{g}' = \mathfrak{g}_{(k')}\)。证明每个 \(\mathfrak{g}'\)-模都由某个 \(\mathfrak{g}\)-模经扩张标量得到，且该模在同构意义下唯一。
\item
  令 \(\mathbf{E}\) 为 \(\mathfrak{g}\)-模。证明下列条件等价：
\end{enumerate}

\begin{enumerate}
\def\labelenumi{(\roman{enumi})}
\item
  E 是忠实的（即~从 g 到 gl(E) 的典则映射是单射）。
\item
  g 的每个根都是 E 的两个权之差。
\end{enumerate}

¶6) 令 \(\varphi\) 为 \(\mathfrak{g}\) 的一个对合自同构，其在 \(\mathfrak{h}\) 上的限制为 -Id。证明若 E 是 \(\mathfrak{g}\)-模，则 E 上存在非退化对称双线性型 \(\varPsi\)，使得

\[
\varPsi (x. a, b) + \varPsi (a, \varphi (x). b) = 0 \quad \text {   for   } x \in \mathfrak {g}, a, b \in \mathrm{E}.
\]

（化归到 E 为单的情形。证明 E 经 \(\varphi\) 变换所得的模同构于 E 的对偶 \(E^{*}\)，从而推出存在满足上述条件的非退化双线性型 \(\Psi\)。接着证明若 e 是 E 中的本原元，则 \(\Psi(e,e)\neq0\)。推出 \(\Psi\) 是对称的。）

\begin{enumerate}
\def\labelenumi{\arabic{enumi})}
\setcounter{enumi}{6}
\tightlist
\item
  若 \(\lambda \in \mathrm{P}_{++}\)，以 \(\rho_{\lambda}:\mathrm{U}(\mathfrak{g})\to \mathrm{End}(\mathrm{E}(\lambda))\) 表示由单模 \(\mathrm{E}(\lambda)\) 定义的表示。我们有 \(\mathrm{Im}(\rho_{\lambda}) = \mathrm{End}(\mathrm{E}(\lambda));\) 置 \(\mathfrak{m}_{\lambda} = \mathrm{Ker}(\rho_{\lambda})\)。\(a)\) 证明 \(\mathfrak{m}_{\lambda}\) 两两不同，并且它们正是这样的所有双边理想：\(\mathfrak{m}\) 属于 \(\mathrm{U}(\mathfrak{g})\)，且 \(\mathrm{U}(\mathfrak{g}) / \mathfrak{m}\) 为有限维单 \(k\)-代数。
\end{enumerate}

\begin{enumerate}
\def\labelenumi{\alph{enumi})}
\setcounter{enumi}{1}
\item
  若 I 是 \(P_{++}\) 的有限子集，置 \(m_{I} = \bigcap_{\lambda \in I} m_{\lambda}\)。证明典则映射 \(U(\mathfrak{g}) / \mathfrak{m}_{I} \to \prod_{\lambda \in I} U(\mathfrak{g}) / \mathfrak{m}_{\lambda}\) 是同构，并且 \(U(\mathfrak{g})\) 的每个有限余维双边理想恰好等于某个 \(m_{I}\)。
\item
  证明 \(\mathrm{U}(\mathfrak{g})\) 的主反自同构将 \(\mathfrak{m}_{\lambda}\) 变为 \(\mathfrak{m}_{\lambda^*}\)，其中 \(\lambda^{*} = -w_{0}\lambda\)（参见~命题11）。
\end{enumerate}

¶8) 令 g 为半单李代数；特别地，在本习题中不假定 g 是分裂的。令 \(\bar{k}\) 为 k 的代数闭包，并令

\[
\pi : \operatorname{Gal} (\bar {k} / k) \to \operatorname{Aut} (\bar {\mathrm{R}}, \bar {\mathrm{B}})
\]

为§5习题8中定义的同态。令 \(\operatorname{Gal}(\bar{k}/k)\) 通过 \(\pi\) 作用于 \(\bar{P}_{++}\)（\(\bar{R}\) 的相对于 \(\bar{B}\) 的支配权集合）；令 \(\Omega\) 为商 \(\bar{P}_{++}/\operatorname{Gal}(\bar{k}/k)\) 中元素的代表元系。

\begin{enumerate}
\def\labelenumi{\alph{enumi})}
\item
  置 \(\bar{g} = \bar{k} \otimes_{k} g\)。若 I 是 \(\bar{P}_{++}\) 的在 \(\operatorname{Gal}(\bar{k}/k)\) 下稳定的有限子集，则与 I 相关的双边理想 \(\bar{m}_{I}\) 是 \(\operatorname{U}(\bar{g})\) 的双边理想，形如 \(\bar{k} \otimes_{k} m_{I}\)，其中 \(m_{I}\) 是 \(\operatorname{U}(g)\) 的双边理想。证明 \(\operatorname{U}(g)\) 的每个有限余维双边理想都唯一地以这种方式得到。
\item
  令 \(\omega \in \Omega\)，令 \(\mathrm{I}(\omega)\) 为它在 \(\operatorname{Gal}(\bar{k}/k)\) 下的轨道，令 \(\mathrm{G}_{\omega}\) 为 \(\omega\) 的稳定子；令 \(k_{\omega}\) 为 \(\bar{k}\) 的子扩张，按 Galois 理论与 \(\mathrm{G}_{\omega}\) 对应。证明 \(\mathrm{U}(\mathfrak{g})/\mathfrak{m}_{\mathrm{I}(\omega)}\) 是中心同构于 \(k_{\omega}\) 的单代数。每个双边理想 \(\mathfrak{m}\)（\(\mathrm{U}(\mathfrak{g})\) 中的），若 \(\mathrm{U}(\mathfrak{g})/\mathfrak{m}\) 是有限维单代数，则  恰好是某个 \(\mathfrak{m}_{\mathrm{I}(\omega)}\)。
\item
  群 \(\operatorname{Gal}(\bar{k}/k)\) 通过 \(\pi\) 作用于表示环 \(\operatorname{R}(\mathfrak{g})\)；以 \(\operatorname{R}(\mathfrak{g})^{\mathrm{inv}}\) 表示 \(\operatorname{R}(\mathfrak{g})\) 中由 \(\operatorname{Gal}(\bar{k}/k)\) 不变元素组成的子环。证明映射 \([E] \to [\bar{k} \otimes_{k} E]\) 延拓为从 \(\operatorname{R}(\mathfrak{g})\) 到 \(\operatorname{R}(\mathfrak{g})^{\mathrm{inv}}\) 的单射同态，且其余核是挠群；该同态是同构，当且仅当对所有 \(\omega \in \Omega\)，\(\operatorname{U}(\mathfrak{g})/\mathfrak{m}_{\operatorname{I}(\omega)}\) 都是 \(k_{\omega}\) 上的矩阵代数；证明当 g 有 Borel 子代数时确实如此。\(^{14}\)
\end{enumerate}

\begin{enumerate}
\def\labelenumi{\arabic{enumi})}
\setcounter{enumi}{8}
\tightlist
\item
  令 \(\mathfrak{a}_1\) 和 \(\mathfrak{a}_2\) 为李代数。证明存在唯一的同态
\end{enumerate}

\[
f: \mathrm{R} (\mathfrak {a} _ {1}) \otimes_ {\mathbf {Z}} \mathrm{R} (\mathfrak {a} _ {2}) \to \mathrm{R} (\mathfrak {a} _ {1} \times \mathfrak {a} _ {2})
\]

使得 \(f([\mathrm{E}_1] \otimes [\mathrm{E}_2]) = [\mathrm{E}_1 \otimes \mathrm{E}_2]\)，当 \(\mathrm{E}_i\) 是 \(\mathfrak{a}_i\)-模（\(i = 1, 2\)）时。证明 \(f\) 是单射，并且当 \(\mathfrak{a}_1\) 和 \(\mathfrak{a}_2\) 为可分裂半单时为双射。

\begin{enumerate}
\def\labelenumi{\arabic{enumi})}
\setcounter{enumi}{9}
\tightlist
\item
  令 \(\Gamma\) 为 \(P\) 的包含 \(Q\) 的子群。这样的子群在 \(W\) 下稳定。
\end{enumerate}

\begin{enumerate}
\def\labelenumi{\alph{enumi})}
\item
  证明若 \(\lambda \in \mathrm{P}_{++} \cap \Gamma\)，则 \(\operatorname{E}(\lambda)\) 的每个权都属于 \(\Gamma\)。
\item
  令 \(\mathrm{R}_{\Gamma}(\mathfrak{g})\) 为 \(\mathrm{R}(\mathfrak{g})\) 的子群，其基为 \([\lambda]\)，其中 \(\lambda \in \mathrm{P}_{++} \cap \Gamma\)。若 \(\mathrm{E}\) 是 \(\mathfrak{g}\)-模，则 \([\mathrm{E}] \in \mathrm{R}_{\Gamma}(\mathfrak{g})\) 当且仅当 \(\mathrm{E}\) 的权属于 \(\Gamma\)。推出 \(\mathrm{R}_{\Gamma}(\mathfrak{g})\) 是 \(\mathrm{R}(\mathfrak{g})\) 的子环。
\item
  证明同态 \(\operatorname{ch} : \mathrm{R}_{\Gamma}(\mathfrak{g}) \to \mathbf{Z}[\Gamma]\) 是从 \(\mathrm{R}_{\Gamma}(\mathfrak{g})\) 到 \(\mathbf{Z}[\Gamma]\) 中由 W 不变元素组成的子环的同构。（使用定理2 (ii)。）
\end{enumerate}

\(d\) ) 具体描述 \(\mathrm{R}(\mathfrak{g})\) 和 \(\mathrm{R}_{\Gamma}(\mathfrak{g})\)：当 \(\mathfrak{g} = \mathfrak{sl}(2,k)\) 且 \(\Gamma = \mathrm{Q}\) 时。

¶11) 记号同第7节。对任意整数 \(m \geq 1\)，以 \(\Psi^{m}\) 表示 \(\mathbf{Z}[\Delta]\) 的自同态，它将 \(e^{\lambda}\) 变为 \(e^{m\lambda}\)。有 \(\Psi^{1} = \mathrm{Id}\) 且 \(\Psi^{m} \circ \Psi^{n} = \Psi^{mn}\)。

令 E 为有限维的 \(\Delta\)-分次向量空间。对所有 \(n \geq 0\)，以 \(a_{n}E\)（分别地，\(s_{n}E\)）表示 E 的第 n 个外幂（分别地，对称幂），并赋予其自然分次。

\begin{enumerate}
\def\labelenumi{\alph{enumi})}
\tightlist
\item
  证明
\end{enumerate}

\[
n \operatorname{ch} (s _ {n} \mathrm{E}) = \sum_ {m = 1} ^ {n} \Psi^ {m} (\operatorname{ch} (\mathrm{E})) \operatorname{ch} (s _ {n - m} \mathrm{E})
\]

以及

\[
n \operatorname{ch} (a _ {n} \mathrm{E}) = \sum_ {m = 1} ^ {n} (- 1) ^ {m - 1} \varPsi^ {m} (\operatorname{ch} (\mathrm{E})) \operatorname{ch} (a _ {n - m} \mathrm{E}).
\]

推出 \(\operatorname{ch}(s_{n}\operatorname{E})\) 和 \(\operatorname{ch}(a_{n}\operatorname{E})\) 可以表示为 \(\varPsi^{m}(\operatorname{ch}(\operatorname{E}))\)（\(1 \leq m \leq n\)）的有理系数多项式。例如：

\[
\mathrm{ch} (s _ {2} \mathrm{E}) = \frac {1}{2} \mathrm{ch} (\mathrm{E}) ^ {2} + \frac {1}{2} \varPsi^ {2} (\mathrm{ch} (\mathrm{E}))
\]

\[
\mathrm{ch} (a _ {2} \mathrm{E}) = \frac {1}{2} \mathrm{ch} (\mathrm{E}) ^ {2} - \frac {1}{2} \varPsi^ {2} (\mathrm{ch} (\mathrm{E})).
\]

\begin{enumerate}
\def\labelenumi{\alph{enumi})}
\setcounter{enumi}{1}
\tightlist
\item
  证明下列恒等式（在变量 T 的形式幂级数代数中，系数属于 \(Q[\Delta]\)）：
\end{enumerate}

\[
\sum_ {n = 0} ^ {\infty} \mathrm{ch} (s _ {n} \mathrm{E}) \mathrm{T} ^ {n} = \exp \left\{\sum_ {m = 1} ^ {\infty} \varPsi^ {m} (\mathrm{ch(E)}) \mathrm{T} ^ {m} / m \right\}
\]

以及

\[
\sum_ {n = 0} ^ {\infty} \mathrm{ch} (a _ {n} \mathrm{E}) \mathrm{T} ^ {n} = \exp \left\{\sum_ {m = 1} ^ {\infty} (- 1) ^ {m - 1} \varPsi^ {m} (\mathrm{ch(E)}) \mathrm{T} ^ {m} / m \right\}.
\]

\begin{enumerate}
\def\labelenumi{\alph{enumi})}
\setcounter{enumi}{2}
\tightlist
\item
  假设 \(\Delta\) 是一个群，从而 \(\varPsi^{m}\) 对所有 \(m\in \mathbf{Z}\) 都可以定义。证明若 \(\mathrm{E}^*\) 是 \(\mathrm{E}\) 的分次对偶，则
\end{enumerate}

\[
\operatorname{ch} \left(\mathrm{E} ^ {*}\right) = \Psi^ {- 1} (\operatorname{ch} (\mathrm{E})).
\]

\begin{enumerate}
\def\labelenumi{\alph{enumi})}
\setcounter{enumi}{3}
\tightlist
\item
  借助 ch 将 \(\mathrm{R}(\mathfrak{g})\) 与 Z[P] 的一个子环相识别。证明 \(\mathrm{R}(\mathfrak{g})\) 在 \(\varPsi^{m}\) 下稳定（\(m \in Z\)），并且前一习题中定义的子环 \(\mathrm{R}_{\varGamma}(\mathfrak{g})\) 也具有此性质。
\end{enumerate}

\begin{enumerate}
\def\labelenumi{\arabic{enumi})}
\setcounter{enumi}{11}
\tightlist
\item
  令 \(\lambda \in \mathrm{P}_{++}\)，令 \(\mathcal{X}_{\lambda}\) 为 \(\operatorname{E}(\lambda)\) 的权集。证明
\end{enumerate}

\[
\mathcal {X} _ {\lambda} \subset \lambda - \mathrm{P} _ {+ +}
\]

一般并不成立。（例如考虑 \(\mathfrak{sl}(3,k)\) 的伴随表示。）

\begin{enumerate}
\def\labelenumi{\arabic{enumi})}
\setcounter{enumi}{12}
\tightlist
\item
  令 \(\lambda = \sum_{\alpha \in B} a_{\alpha} \alpha\) 为 \(P_{++}\) 中的一个元素。对于 \(n = 0, 1, \ldots\)，以 \(\mathcal{X}_n\) 表示权 \(\mu\)，它是 \(E(\lambda)\) 的权，且满足 \(\lambda - \mu\) 是 n 个 \(n\)\(B\) 中元素之和。令 \(s_n\) 为 \(\mathcal{X}_n\) 中元素的重数之和（作为 \(E(\lambda)\) 的权的重数）。令 \(T = 2 \sum_{\alpha \in B} a_{\alpha}\)。证明：
\end{enumerate}

\begin{enumerate}
\def\labelenumi{\alph{enumi})}
\item
  T 是一个整数 \(\geq 0\)。
\item
  当 n > T 时，\(s_{n}=0\)，且 \(s_{T-n}=s_{n}\)。
\item
  若 r 是 T/2 的整数部分，则 \(s_{1} \leq s_{2} \leq \cdots \leq s_{r+1}\)。
\end{enumerate}

¶14) 令 \(\lambda \in \mathrm{P}_{++}\)，令 \(\mathrm{F}_{\lambda}\) 为 \(Z(\lambda)\) 的最大真子模，令 \(v\) 为 \(Z(\lambda)\) 中权为 \(\lambda\) 的本原元，参见~§6第3节。证明

\[
\mathrm{F} _ {\lambda} = \sum_ {\alpha \in \mathrm{B}} \mathrm{U} (\mathfrak {g}) X _ {- \alpha} ^ {\lambda (H _ {\alpha}) + 1} v = \sum_ {\alpha \in \mathrm{B}} \mathrm{U} (\mathfrak {n} _ {-}) X _ {- \alpha} ^ {\lambda (H _ {\alpha}) + 1} v.
\]

\begin{enumerate}
\def\labelenumi{\arabic{enumi})}
\setcounter{enumi}{14}
\tightlist
\item
  令 \(\lambda \in \mathfrak{h}^*\)，令 \(v\)（分别地，\(v'\)）为 \(\mathrm{Z}(\lambda)\)（分别地，\(\mathrm{E}(\lambda)\)）中权为 \(\lambda\) 的本原元。令 \(\mathrm{I}\)（分别地，\(\mathrm{I}'\)）为 \(v\)（分别地，\(v'\)）在 \(\mathrm{U}(\mathfrak{g})\) 中的湮没子。
\end{enumerate}

\[
a) \mathrm{I} = \mathrm{U} (\mathfrak {g}) \mathfrak {n} _ {+} + \sum_ {h \in \mathfrak {h}} \mathrm{U} (\mathfrak {g}) (h - \lambda (h)).
\]

\begin{enumerate}
\def\labelenumi{\alph{enumi})}
\setcounter{enumi}{1}
\item
I' 是不等于 U(g) 且包含 I 的 U(g) 的最大左理想。
\item
  若 \(\lambda \in \mathrm{P}_{++}\)，则
\end{enumerate}

\[
\mathrm{I} ^ {\prime} = \mathrm{I} + \sum_ {\alpha \in \mathrm{B}} \mathrm{U} (\mathfrak {g}) X _ {- \alpha} ^ {\lambda (H _ {\alpha}) + 1} = \mathrm{I} + \sum_ {\alpha \in \mathrm{B}} \mathrm{U} (\mathfrak {n} _ {-}) X _ {- \alpha} ^ {\lambda (H _ {\alpha}) + 1}.
\]

（使用前一习题。）

¶16) 令 V 和 V' 为两个 g-模。若存在一个线性映射 \(f : V \to V'\) 使其满足下列条件，则称 V' 从属于 V：

\(\alpha)\) \(f\) 是满射；\(\beta)\) \(f\) 下，\(\mathbf{V}\) 的本原元的像或者为 0，或者是 \(\mathrm{V}^{\prime}\) 的本原元；\(\gamma)\) \(f\) 是 \(\mathfrak{n}_{-}\)-同态。

\begin{enumerate}
\def\labelenumi{\alph{enumi})}
\item
  令 \(f: V \to V'\) 满足 \(\alpha\) )、\(\beta\) )、\(\gamma\) )。令 v 为 V 的本原元。则由 v 生成的 g-子模在 f 下的像是由 \(f(v)\) 生成的 g-子模。
\item
  令 \(f: \mathrm{V} \to \mathrm{V}'\) 满足 \(\alpha), \beta), \gamma)\)。令 \(\mathrm{W}\) 为 \(\mathfrak{g}\) 的 \(\mathrm{V}\)-子模。则 \(f(\mathrm{W})\) 是 \(\mathfrak{g}\) 的 \(\mathrm{V}'\)-子模，并且从属于 \(\mathrm{W}\)。
  \item
  令 \(V = E_{1} \oplus \cdots \oplus E_{s}\) 且 \(V' = E_{1}' \oplus \cdots E_{s'}\) 为 V 和 \(V'\) 分解为单模之和的分解。则 \(V'\) 从属于 V，当且仅当 \(s' \leq s\)，且存在 \(\sigma \in S_{s}\)，使得 \(E_{i}'\) 从属于 \(\mathrm{E}_{\sigma(i)}\)，对 \(i = 1, \ldots, s'\) 成立。
\item
  若 \(V'\) 从属于 V 且 V 是单的，则 \(V'\) 是单的或化为 0。
\item
  假设 V 和 \(V'\) 是单的。令 \(\lambda\) 和 \(\lambda'\) 为它们的最高权。则 \(V'\) 从属于 V，当且仅当 \(\lambda'(H_{\alpha}) \leq \lambda(H_{\alpha})\) 对所有 \(\alpha \in B\) 成立。（充分性使用前一习题。）
\end{enumerate}

\begin{enumerate}
\def\labelenumi{\arabic{enumi})}
\setcounter{enumi}{16}
\tightlist
\item
  令 \(\lambda, \mu \in \mathrm{P}_{++}\) 和 \(\alpha \in \mathrm{B}\) 满足 \(\lambda(H_{\alpha}) \geq 1\) 且 \(\mu(H_{\alpha}) \geq 1\)。令 \(\mathrm{F} = \mathrm{E}(\lambda) \otimes \mathrm{E}(\mu)\)。
\end{enumerate}

\begin{enumerate}
\def\labelenumi{\alph{enumi})}
\item
  证明 \(\dim \mathrm{F}^{\lambda +\mu} = 1\)，且 \(\dim \mathrm{F}^{\lambda +\mu -\alpha} = 2\)。
\item
  证明 \(X_{\alpha}: F^{\lambda+\mu-\alpha} \to F^{\lambda+\mu}\) 是满射，且其核中的非零元素是本原元（注意，若 \(\beta \in B\) 异于 \(\alpha\)，则 \(\lambda + \mu - \alpha + \beta\) 不是 F 的权）。
\item
  推出 \(\mathrm{E}(\lambda)\otimes\mathrm{E}(\mu)\) 含有唯一一个同构于
\end{enumerate}

\[
\mathrm{E} (\lambda + \mu - \alpha).
\]

\begin{enumerate}
\def\labelenumi{\alph{enumi})}
\setcounter{enumi}{3}
\tightlist
\item
  证明 \(\mathbf{S}^{2}(\mathrm{E}(\lambda))\)（分别地，\(\bigwedge^{2}\mathrm{E}(\lambda)\)）含有唯一一个同构于 \(\mathrm{E}(2\lambda)\)（分别地，同构于 \(\mathrm{E}(2\lambda-\alpha)\)）的子模。
\end{enumerate}

¶ 18) 选取一个正定非退化对称双线性形式 \((\cdot|\cdot)\) 在 \(\mathfrak{h}_{\mathbf{R}}^{*}\) 上 W-不变。令 \(\lambda \in P_{++}\)。

\begin{enumerate}
\def\labelenumi{\alph{enumi})}
\item
  令 \(\mu\) 为 \(\mathrm{E}(\lambda)\) 的一个权。将 \(\lambda - \mu\) 写成 \(\sum_{\alpha \in \mathrm{B}} k_{\alpha} \alpha\) 的形式。令 \(\alpha \in \mathrm{B}\) 满足 \(k_{\alpha} \neq 0\)。证明存在 \(\alpha_{1}, \ldots, \alpha_{n} \in \mathrm{B}\)，使得 \((\lambda | \alpha_{1}) \neq 0\)、\((\alpha_{1} | \alpha_{2}) \neq 0, \ldots, (\alpha_{n-1} | \alpha_{n}) \neq 0\)、\((\alpha_{n} | \alpha) \neq 0\)。
\item
  令 \(v\) 为 \(\mathrm{E}(\lambda)\) 的本原元，并令 \(\alpha_1, \ldots, \alpha_n \in \mathrm{B}\) 满足下列条件：
\end{enumerate}

\begin{enumerate}
\def\labelenumi{(\roman{enumi})}
\item
  \((\alpha_{i}|\alpha_{i+1})\neq0\) 对 \(i=1,2,\ldots,n-1;\)
\item
  \((\alpha_{i}|\alpha_{j}) = 0\) 对 \(j > i + 1\)；
\item
  \(\lambda(H_{\alpha_{1}})\neq0\) 且 \(\lambda(H_{\alpha_{2}})=\cdots=\lambda(H_{\alpha_{n}})=0.\)
\end{enumerate}

证明 \(X_{-\alpha_n}X_{-\alpha_{n-1}} \ldots X_{-\alpha_1}v \neq 0\)。(注意，对于 \(1 \leq s \leq n\)，

\[
\lambda - \alpha_ {1} - \dots - \alpha_ {s - 1} + \alpha_ {n}
\]

不是 \(\mathrm{E}(\lambda)\) 的权，并由归纳法（以 \(X_{\alpha_s}X_{-\alpha_s}X_{-\alpha_{s-1}}\ldots X_{-\alpha_1}v\neq 0\) 为归纳变量）推出 \(s\)。) 若 \(\sigma \in \mathfrak{S}_n\) 且 \(\sigma \neq 1\)，则 \(X_{-\alpha_{\sigma(n)}}X_{-\alpha_{\sigma(n-1)}}\ldots X_{-\alpha_{\sigma(1)}}v = 0\)。(令 \(r\) 为满足 \(\sigma(r)\neq r\) 的最小整数。利用 \(a\) ) 证明 \(\lambda -\alpha_{\sigma(1)} - \ldots -\alpha_{\sigma(r)}\) 不是 \(\mathrm{E}(\lambda)\) 的权。)

\begin{enumerate}
\def\labelenumi{\alph{enumi})}
\setcounter{enumi}{2}
\tightlist
\item
  令 \(\lambda' \in \mathrm{P}_{++}\)。连接 \(\lambda\) 与 \(\lambda'\) 的链是 B 中元素的序列 \((\alpha_1, \ldots, \alpha_n)\)，满足 \(n \geq 1\)、\((\lambda|\alpha_1) \neq 0\)、\((\alpha_1|\alpha_2) \neq 0, \ldots, (\alpha_{n-1}|\alpha_n) \neq 0\)、\((\alpha_n|\lambda') \neq 0\)。若这样的链的 \((\alpha_1, \ldots, \alpha_n)\) 没有连接 \(\lambda\) 与 \(\lambda'\) 的真子序列，则称其为极小链。在此情形，有：\((\alpha_i|\alpha_j) = 0\)（若 \(|i - j| \geq 2\)）；若 \((\lambda|\alpha_i) = 0\)，则 \(i \geq 2\)；若 \((\lambda'|\alpha_i) = 0\)，则 \(i \leq n - 1\)。
\end{enumerate}

\(d\) ) 令 \((\alpha_{1},\ldots ,\alpha_{n})\) 为连接 \(\lambda\) 与 \(\lambda^{\prime}\) 的极小链。若 \(v^{\prime}\) 是 \(\operatorname {E}(\lambda ')\) 中的本原向量，置：

\[
v _ {s} = X _ {- \alpha_ {s}} X _ {- \alpha_ {s - 1}} \ldots X _ {- \alpha_ {1}} v (s = 0, 1, \ldots , n)
\]

\[
v _ {s} ^ {\prime} = X _ {- \alpha_ {s + 1}} X _ {- \alpha_ {s + 2}} \dots X _ {- \alpha_ {n}} v ^ {\prime} (s = 0, 1, \ldots , n)
\]

\(a_0 = (\lambda |\alpha_1), a_n = (-1)^n (\lambda '|\alpha_n), a_s = (-1)^{s + 1}(\alpha_s|\alpha_{s + 1}), 1 \leq s \leq n - 1.\) 证明

\[
\sum_ {s = 0} ^ {n} a _ {s} v _ {s} \otimes v _ {s} ^ {\prime}
\]

是 \(\mathrm{E}(\lambda)\otimes\mathrm{E}(\lambda')\) 中权为 \(\lambda+\lambda'-\alpha_{1}-\cdots-\alpha_{n}\) 的本原元，并且在相差一个伸缩的意义下唯一。

（使用 b) 和 c) 证明 \(\mathrm{E}(\lambda) \otimes \mathrm{E}(\lambda')\) 中每个权为

\[
\lambda + \lambda^ {\prime} - \alpha_ {1} - \dots - \alpha_ {n}
\]

的是 \(v_{s} \otimes v_{s}^{\prime}\) 的线性组合。然后写出这种线性组合成为本原向量的条件。）

推出 \(\mathrm{E}(\lambda) \otimes \mathrm{E}(\lambda')\) 含有唯一一个同构于

\[
\operatorname{E} \left(\lambda + \lambda^ {\prime} - \alpha_ {1} - \dots - \alpha_ {n}\right).
\]

（当 \(n = 1\) 时，我们得到习题17。）

\begin{enumerate}
\def\labelenumi{\alph{enumi})}
\setcounter{enumi}{4}
\tightlist
\item
  令 w 为 \(\mathrm{E}(\lambda)\otimes\mathrm{E}(\lambda')\) 的本原元。假设 w 的权 \(\nu\) 异于 \(\lambda+\lambda'\)。证明存在一个 \((\alpha_{1},\ldots,\alpha_{n})\)，它连接 \(\lambda\) 与 \(\lambda'\)，使得
\end{enumerate}

\[
\nu \leq \lambda + \lambda^ {\prime} - \alpha_ {1} - \dots - \alpha_ {n}.
\]

(令 C 为由 \(\alpha \in \mathrm{B}\) 中满足 \(\alpha\) 在 \(\lambda + \lambda' - \nu\) 中的指标为 \(\neq 0\) 的元素组成的集合。令 D（分别地，D'）为满足下述条件的 \(\alpha \in \mathrm{C}\) 的集合：存在 \(\gamma_1, \ldots, \gamma_t \in \mathrm{C}\)，使得 \((\lambda|\gamma_1) \neq 0\)（分别地，\((\lambda'|\gamma_1) \neq 0\)）、\((\gamma_1|\gamma_2) \neq 0, \ldots, (\gamma_{t-1}|\gamma_t) \neq 0\)、\((\gamma_t|\alpha) \neq 0\)。令 Y（分别地，Y'）为 E( \(\lambda\) )（分别地，E( \(\lambda'\) )）中形如 \(\lambda - \sum_{\alpha \in \mathrm{C}} k_\alpha\alpha\)（分别地，\(\lambda' - \sum_{\alpha \in \mathrm{C}} k_\alpha\alpha\)）且 \(k_\alpha \in \mathbf{N}\) 的权的集合。证明 \(w\) 属于 \(\left(\sum_{\mu \in \mathrm{Y}}\mathrm{E}(\lambda)^{\mu}\right)\otimes \left(\sum_{\mu^{\prime}\in \mathrm{Y}^{\prime}}\mathrm{E}(\lambda^{\prime})^{\mu^{\prime}}\right)\)。利用 \(a\) ) 以及 \(\nu \neq \lambda +\lambda '\) 这一事实，证明 \(\mathrm{D}\cap \mathrm{D}'\neq \emptyset .\))

\(f\) ) 证明下列性质等价：

\begin{enumerate}
\def\labelenumi{(\roman{enumi})}
\item
  \(\mathrm{E}(\lambda)\otimes \mathrm{E}(\lambda^{\prime})\) 同构于 \(\mathrm{E}(\lambda +\lambda^{\prime})\)
\item
  \(\mathrm{E}(\lambda) \otimes \mathrm{E}(\lambda')\) 是单模。
\item
  不存在连接 \(\lambda\) 与 \(\lambda'\) 的链。
\item
  R 是两个根系 \(R_{1}\) 和 \(R_{1}^{\prime}\) 的直和，使得 \(\lambda \in \mathrm{P}(\mathrm{R}_{1})\) 且 \(\lambda^{\prime} \in \mathrm{P}(\mathrm{R}_{1}^{\prime})\)。
\item
  \(\mathfrak{g}\) 是两个理想 \(\mathfrak{s}\) 和 \(\mathfrak{s}'\) 的直积，使得 \(\mathfrak{s}'\) .E(λ) = 0 且 \(\mathfrak{s}\) .E(λ') = 0。
\end{enumerate}

（使用 \(d\) ) 证明 (ii) 与 (iii) 等价。）\(^{15}\)

\begin{enumerate}
\def\labelenumi{\arabic{enumi})}
\setcounter{enumi}{18}
\item
  回忆命题10的记号。令 \(\bar{k}\) 为 \(k\) 的代数闭包。若 \(x \in \mathfrak{g}\)，以 \(\mathcal{X}_{\mathrm{E}}(x)\)（分别地，\(\mathcal{X}_{\mathrm{F}}(x)\)、分别地，\(\mathcal{X}_{\mathrm{G}}(x)\)）表示 \(x_{\mathrm{E}}\)（分别地，\(x_{\mathrm{F}}\)、分别地，\(x_{\mathrm{G}}\)）在 \(\bar{k}\) 中的特征值集合。证明 \(\mathcal{X}_{\mathrm{G}}(x) = \mathcal{X}_{\mathrm{E}}(x) + \mathcal{X}_{\mathrm{F}}(x)\) 对所有 \(x \in \mathfrak{g}\) 成立，并且在给定 E 和 F 时，此性质在同构意义下刻画单 \(\mathfrak{g}\)-模 G。
\item
  假设 \(k = \mathbf{R}\) 或 \(\mathbf{C}\)。令 \(\Gamma\) 为李代数为 \(\mathfrak{g}\) 的单连通李群。令 \(\lambda, \mu, \mathrm{E}, \mathrm{F}, \mathrm{G}\) 如命题10中。令 \((e_1, \ldots, e_n)\)（分别地，\((f_1, \ldots, f_p)\)）为 \(\mathrm{E}\)（分别地，\(\mathrm{F}\)）的一个由 \(\mathfrak{h}\) 的特征向量组成的基，其中 \(e_1 \in \mathrm{E}^\lambda\) 且 \(f_1 \in \mathrm{F}^\mu\)。可将 \(\mathrm{E}, \mathrm{F}, \mathrm{G}\) 看作 \(\Gamma\)-模。若 \(\gamma \in \Gamma\)，以 \(a_i(\gamma)\) 表示 \(\gamma.e_1\) 的指标为 \(i\) 的坐标，以 \(b_j(\gamma)\) 表示 \(\gamma.f_1\) 的指标为 \(j\) 的坐标。证明 \(a_i b_j\) 上的函数 \(\Gamma\) 不恒为零。推出对所有 \((i,j)\)，存在 \(\mathrm{G} \subset \mathrm{E} \otimes \mathrm{F}\) 中的元素，其指标为 \((i,j)\) 的坐标 \(\neq 0\)；当 \(k = \mathbf{R}\) 或 \(\mathbf{C}\) 时，这给出命题10的一个新证明。由此转到 \(k = \mathbf{Q}\) 的情形，再转到任意域的情形，参见~习题4。
\item
  令 \(\lambda, \mu \in \mathrm{P}_{++}\)。令 E、F、G 为单 \(\mathfrak{g}\)-模，其最高权为 \(\lambda, \mu, \lambda + \mu\)，令 \(n\) 为满足 \(\geq 1\) 的整数。若 \(\omega\) 是 E 的一个重数为 \(n\) 的权，则 \(\omega + \mu\) 是 G 的一个重数 \(\geq n\) 的权。（有 \(\mathrm{G}^{\omega + \mu} \subset \bigoplus_{\nu + \sigma = \omega + \mu} \mathrm{E}^{\nu} \otimes \mathrm{F}^{\sigma}\)。若 \(\dim \mathrm{G}^{\omega + \mu} < n\)，则 \(\mathrm{G}^{\omega + \mu}\) 到 \(\mathrm{E}^{\omega} \otimes \mathrm{F}^{\mu}\) 的投影形如 \(\mathrm{E}' \otimes \mathrm{F}^{\mu}\)，其中 \(\mathrm{E}'\) 真包含于 \(\mathrm{E}^{\omega}\)。通过选取 E 和 F 的适配基并仿照命题10的证明，由此导出矛盾。）
\item
  假设 \(\mathfrak{g}\) 是单的，换言之，R 是不可约的。证明存在唯一的支配权 \(\lambda \neq 0\)，使得 \(\mathrm{E}(\lambda)\) 的权集为 \(\mathrm{W}.\lambda \cup \{0\}\)：有 \(\lambda = \alpha\)，其中 \(\alpha \in \mathrm{R}\) 满足 \(H_{\alpha}\) 是 \(\mathrm{R}^{\vee}\) 的最高根。当所有根等长时（情形 \(\mathrm{A}_l,\mathrm{D}_l,\mathrm{E}_6,\mathrm{E}_7,\mathrm{E}_8\)），有 \(\lambda = \tilde{\alpha}\)；这是唯一一个同时为支配权的根；相应的表示是 g 的伴随表示。在其他情形中，\(\lambda\) 是长度最小且为支配权的唯一根；按第VI章图版中的记号，有：\(\lambda = \varpi_{1}\)（\(B_{l}\) 型）；\(\lambda = \varpi_{2}\)（\(C_{l}\) 型）；\(\lambda = \varpi_{4}\)（\(F_{4}\) 型）；\(\lambda = \varpi_{1}\)（\(G_{2}\) 型）。
\item
  令 \(\mathrm{U}^0\) 为 \(\mathfrak{h}\) 在 \(\mathrm{U}(\mathfrak{g})\) 中的交换子（参见~§6第4节）。
\end{enumerate}

\begin{enumerate}
\def\labelenumi{\alph{enumi})}
\item
  令 V 为一个（有限维）g-模。证明各 \(V^{\lambda}\)（\(\lambda \in P\)）在 \(U^{0}\) 下稳定，并且若 V 是单的且 \(V^{\lambda} \neq 0\)，则 \(V^{\lambda}\) 是单 \(U^{0}\)-模（使用 \(\mathrm{U}(\mathfrak{g}) = \oplus\mathrm{U}^{\lambda}\) 的分解，同上）。
\item
  证明对于每个 \(c \neq 0\)（属于 \(U^{0}\)），都存在 \(\rho\) 的有限维单表示 \(U^{0}\)，使得 \(\rho(c) \neq 0\)。（使用 a) 和第I章§7的习题3 a)。
\end{enumerate}

\begin{enumerate}
\def\labelenumi{\arabic{enumi})}
\setcounter{enumi}{23}
\tightlist
\item
  令 \(A\) 为幺元结合代数，\(M\) 为其有限余维双边理想的集合。
\end{enumerate}

\begin{enumerate}
\def\labelenumi{\alph{enumi})}
\tightlist
\item
  令 \(\mathrm{U}^*\) 为 U 的向量空间对偶。若 \(\theta \in \mathrm{U}^*\)，证明下列性质等价：
\end{enumerate}

\begin{enumerate}
\def\labelenumi{(\roman{enumi})}
\item
  存在 \(\mathfrak{m} \in \mathrm{M}\)，使得 \(\theta(\mathfrak{m}) = 0\)；
\item
   存在两个有限族 \((\theta_{i}^{\prime})\) 和 \((\theta_{i}^{\prime\prime})\)，它们由 \(U^{*}\) 中的元素组成，使得
\end{enumerate}

\[
\theta (x y) = \sum_ {i} \theta_ {i} ^ {\prime} (x) \theta_ {i} ^ {\prime \prime} (y) \quad \text { for   all } x, y \in \mathrm{U}.
\]

具有这些性质的元素 \(\theta\) 构成 \(\mathrm{U}'\) 的子空间，它与《代数》第III章§11习题27中记为 \(\mathrm{U}^*\)\(\mathrm{B}'\) 的空间重合。\(\mathrm{U}'\) 上存在唯一的余代数结构，其余积 \(c: \mathrm{U}' \to \mathrm{U}' \otimes \mathrm{U}'\) 由下式给出

\[
c (\theta) = \sum_ {i} \theta_ {i} ^ {\prime} \otimes \theta_ {i} ^ {\prime \prime},
\]

其中 \(\theta_i', \theta_i'' \in \mathrm{U}'\) 满足 \(\theta(xy) = \sum_{i} \theta_i'(x) \theta_i''(y)\) 对所有 \(x, y \in \mathrm{U}\) 成立，参见~(ii)。

余代数 \(\mathrm{U}'\) 是由子空间 \((\mathrm{U} / \mathfrak{m})^*\)（\(\mathfrak{m} \in \mathrm{M}\)）构成的递增滤过的并；这些子空间都是有限维的。\(\mathrm{U}'\) 的对偶可以与代数 \(\hat{\mathrm{U}} = \varprojlim \mathrm{U} / \mathfrak{m}\) 相识别；若给 \(\hat{\mathrm{U}}\) 赋予 \(\mathrm{U} / \mathfrak{m}, \mathfrak{m} \in \mathrm{M}\) 上的离散拓扑的射影极限拓扑，则 \(\hat{\mathrm{U}}\) 上的连续线性形式正是 \(\mathrm{U}'\) 中的元素。


\begin{enumerate}
\def\labelenumi{\alph{enumi})}
\setcounter{enumi}{1}
\tightlist
\item
  令 E 为有限维左 U-模。其湮没子 \(m_{E}\) 属于 M；复合 \(\hat{U} \to U/m_{E} \to \text{End}(E)\) 赋予 E 一个左 \(\hat{U}\)-模结构。若 F 为有限维左 U-模，则线性映射 \(f : E \to F\) 是 U-同态，当且仅当它是 \(\hat{U}\)-同态。若 \(a \in E, b \in E^{*}\)，线性形式 \(\theta_{a,b} : x \mapsto \langle xa, b \rangle\) 属于 \(U'\)，且
\end{enumerate}

\[
\langle x, \theta_ {a, b} \rangle = \langle x a, b \rangle \quad \text { for   all } x \in \hat {\mathrm{U}}.
\]

  当 E 变化时，\(\theta_{a,b}\)（其中 \(a,b\) 变化）生成 \(k\)-向量空间 U'。

\begin{enumerate}
\def\labelenumi{\alph{enumi})}
\setcounter{enumi}{2}
\tightlist
\item
  令 \(X_{U}\) 为有限维左 U-模的同构类集合。对所有 \(E \in X_{U}\)，令 \(u_{E}\) 为 E 的 k-线性自同态；假设
\end{enumerate}

\[
f \circ u _ {\mathrm{E}} = u _ {\mathrm{F}} \circ f \quad \text { for   all   E, F } \in X_{U} \text { and } f \in \operatorname{Hom} _ {\mathrm{U}} (\mathrm{E,F}).
\]

证明存在唯一的元素 \(x \in \hat{U}\)，使得 \(x_{E} = u_{E}\) 对所有 \(E \in X_{U}\) 成立。（化归到 U 为有限维的情形。）

¶ 25) 令 a 为李代数，U 为其包络代数。将习题24的定义和结果应用于代数 U。特别地，\(U' \subset U^*\)，且 \(U'\) 的对偶可以与代数 \(\hat{U} = \varprojlim U/m\) 相识别；典则映射 \(U \to \hat{U}\) 是单射（第I章§7的习题3）。

\begin{enumerate}
\def\labelenumi{\alph{enumi})}
\tightlist
\item
  U 的余代数结构（第II章§1第4节）在对偶 U* 上定义一个代数结构（参见~第II章§1第5节的命题10以及《代数》第III章§11第2节）。若 E、F 是有限维 a-模，
\end{enumerate}

\[
\theta_ {a, b} \circ \theta_ {c, d} = \theta_ {a \otimes c, b \otimes d} \text { for } a \in \mathrm{E}, b \in \mathrm{E} ^ {*}, c \in \mathrm{F}, d \in \mathrm{F} ^ {*}.
\]

推出 \(U'\) 是 \(U^{*}\) 的子代数。\(U'\) 的余代数结构和代数结构使其成为交换双代数（《代数》第III章§11第4节）。

\begin{enumerate}
\def\labelenumi{\alph{enumi})}
\setcounter{enumi}{1}
\tightlist
\item
  令 \(x\) 为 \(\hat{\mathrm{U}}\) 的一个元素；将 \(x\) 与线性形式 \(\mathrm{U}' \to k\) 相识别。证明下列性质等价：
\end{enumerate}

\begin{enumerate}
\def\labelenumi{(\roman{enumi})}
\item
  \(x\) 是从 \(\mathrm{U}'\) 到 \(k\) 的代数同态。
\item
  对所有有限维 a-模 E 和 F，都有 \(x_{E \otimes F} = x_{E} \otimes x_{F}\)。
\end{enumerate}

（先证明 (ii) 等价于

\[
\left(\mathrm{ii} ^ {\prime}\right) x \left(\theta_ {a \otimes c, b \otimes d}\right) = x \left(\theta_ {a, b}\right) x \left(\theta_ {c, d}\right) \text {if} a \in \mathrm{E}, b \in \mathrm{E} ^ {*}, c \in \mathrm{F}, d \in \mathrm{F} ^ {*},
\]

并使用 \(\theta_{a,b}\) 生成 \(\mathrm{U}'\) 这一事实。）

\(c)\) 令 \(x\) 为满足 \(\hat{\mathbf{U}}\) ) 中条件 (i) 和 (ii) 的 \(b\) 元素。证明下列条件等价：

\begin{enumerate}
\def\labelenumi{(\roman{enumi})}
\setcounter{enumi}{2}
\item
  \(x\) 将 \(\mathrm{U}'\) 的单位元映为 \(k\) 的单位元。
\item
  \(x \neq 0.\)
\item
  若给 \(k\) 赋予平凡的 \(\mathfrak{a}\)-模结构，则有 \(x_{k} = \mathrm{Id}\)。
\item
  对所有 E，\(x_{\mathrm{E}}\) 都是可逆的。
\end{enumerate}

（(iii) \(\Longleftrightarrow\) (iv) 的等价性由 \(x\) 是代数同态这一事实得到。另一方面，\(x_{k} = \lambda \mathrm{Id}\)，其中 \(\lambda \in k\)。利用 \(\mathfrak{a}\)-同构 \(k \otimes \mathrm{E} \to \mathrm{E}\)，推出对所有 E 都有 \(\lambda x_{\mathrm{E}} = x_{\mathrm{E}}\)，特别地，取 \(\lambda^2 = \lambda\) 得 \(\mathrm{E} = k\)。\(\lambda = 1\) 的情形对应于 \(x \neq 0\)，因而 (iv) \(\Longleftrightarrow\) (v)，且 (vi) \(\Longrightarrow\) (v)。为证明 (v) \(\Longrightarrow\) (vi)，证明若 F 是 E 的对偶，则 \(^t x_{\mathrm{F}} \circ x_{\mathrm{E}} = \lambda \mathrm{Id}_{\mathrm{E}}\)。）

\begin{enumerate}
\def\labelenumi{\alph{enumi})}
\setcounter{enumi}{3}
\tightlist
\item
  令 G 为 \(\hat{U}\) 中满足上述条件 (i) 至 (vi) 的元素集。证明 G 是 \(\hat{U}\) 的可逆元素群的子群。
\end{enumerate}

令 \(x \in G\)。若 E 是有限维 \(\mathfrak{a}\)-模，则 \(x_{\mathrm{E}} \in \mathbf{GL}(\mathrm{E})\)。特别地，取 \(\mathrm{E} = \mathfrak{a}\)，赋予其伴随表示；这给出元素 \(x_{\mathfrak{a}} \in \mathbf{GL}(\mathfrak{a})\)。证明 \(x_{\mathfrak{a}}\) 是 \(\mathfrak{a}\) 的自同构（使用由括号给出的 \(\mathfrak{a}\)-同态 \(\mathfrak{a} \otimes \mathfrak{a} \to \mathfrak{a}\)）。这给出同态 \(v: \mathrm{G} \to \operatorname{Aut}(\mathfrak{a})\)。若 \(\mathrm{E}\) 是有限维 \(\mathfrak{a}\)-模，

\[
x _ {\mathrm{E}} (y. e) = v (x) (y). x _ {\mathrm{E}} (e) \quad \mathrm{if} y \in \mathfrak {a}, e \in \mathrm{E}
\]

（使用由 \(\mathfrak{a}\) 在 \(\mathfrak{a} \otimes \mathrm{E} \to \mathrm{E}\) 上的作用给出的 \(\mathfrak{a}\)-同态 \(\mathrm{E}\)。）

\begin{enumerate}
\def\labelenumi{\alph{enumi})}
\setcounter{enumi}{4}
\tightlist
\item
  U 的主反自同构 \(\sigma\) 由连续性延拓到 \(\hat{\mathbf{U}}\)。其转置保持 \(\mathrm{U}'\) 稳定，并在 \(\mathrm{U}'\) 上诱导逆（《代数》第III章§11习题4）。若 \(x \in \mathrm{G}\)，则 \(\sigma(x) = x^{-1}\)。
\end{enumerate}

\(f\) ) 假设 \(\mathfrak{a}\) 是半单的 \(^{16}\)。令 \(n\) 为 \(\mathfrak{a}\) 的幂零元。则存在唯一的 \(e^n\) 中的元素 \(\mathrm{G}\)，使得对每个有限维 \((e^n)_{\mathrm{E}} = \exp (n_{\mathrm{E}})\)-模 E 都有 \(\mathfrak{a}\)。有 \(v(e^n) = \exp (\operatorname {ad}n)\in \operatorname {Aut}(\mathfrak{a})\)，所以 \(\operatorname {Aut}_e(\mathfrak{a})\subset v(\mathrm{G})\)。

证明若 b 是由幂零元素组成的 a 的子代数，则

\[
e ^ {n}. e ^ {m} = e ^ {\mathrm{H} (n, m)} \quad \text { for } n, m \in \mathfrak {b},
\]

其中 H 表示 Hausdorff 级数（第II章§6）。

¶ 26) 将习题25的记号和结果应用于 a = g 的情形（分裂情形）。

\begin{enumerate}
\def\labelenumi{\alph{enumi})}
\item
  令 \(x \in G\)，令 \(\sigma = v(x)\) 为它在 \(\operatorname{Aut}(\mathfrak{g})\) 中的像。若 \(\rho\) 是 g 的一个表示，则 \(\rho\) 与 \(\rho \circ \sigma\) 等价。推出（参见~第2节，注1）\(\sigma\) 属于 \(\operatorname{Aut}_{0}(\mathfrak{g})\)。将此结果推广到任意半单代数。
\item
  令 \(\varphi \in \mathrm{T}_{\mathrm{P}} = \mathrm{Hom}(\mathrm{P}, k^{*})\)，其中 \(\mathrm{P} = \mathrm{P}(\mathrm{R})\)。若 \(\mathrm{E}\) 是 \(\mathfrak{g}\)-模，令 \(\varphi_{\mathrm{E}}\) 为 \(\mathrm{E}\) 的自同态，其在每个 \(\mathrm{E}^{\lambda}\)（\(\lambda \in \mathrm{P}\)）上的限制是比值为 \(\varphi(\lambda)\) 的伸缩。证明存在唯一的元素 \(t(\varphi) \in \mathrm{G}\)，使得对所有 \(t(\varphi)_{\mathrm{E}} = \varphi_{\mathrm{E}}\) 都有 \(\mathrm{E}\)（使用习题24 c) 以及习题25的刻画 (ii) 和 (vi)）。这给出同态 \(t: \mathrm{T}_{\mathrm{P}} \to \mathrm{G}\)。证明 \(t\) 是单射。利用这一点将 \(\mathrm{T}_{\mathrm{P}}\) 与 \(\mathrm{G}\) 的一个子群相识别。复合映射 \(\mathrm{T}_{\mathrm{P}} \to \mathrm{G} \to \mathrm{Aut}(\mathfrak{g})\) 是§5第2节中记为 \(f \circ q\) 的同态。
\item
  令 \(x \in G\) 满足 \(\sigma = v(x)\) 属于 \(f(\mathrm{T}_{\mathrm{Q}})\) 的子群 \(\mathrm{Aut}_{0}(\mathfrak{g})\)（§5第2节），换言之，它在 h 上作用平凡；以 \(T_{Q} = \mathrm{Hom}(Q, k^{*})\) 表示与 x 对应的 \(\psi\) 中的元素。我们将证明 x 属于 \(T_{P}\)。依次证明：
\end{enumerate}

\(c_{1}\) ) 若 E 是 \(\mathfrak{g}\)-模，则 \(x_{\mathrm{E}}\) 是 E 的 \(\mathfrak{h}\)-自同态。

（使用 \(g \otimes E \to E\) 的 g-同态，以及 x 在 h 上作用平凡这一事实。）特别地，各 \(E^{\mu}\) 在 \(x_{E}\) 下稳定。

\(c_{2}\) ) 存在 \(\varphi \in \mathrm{T}_{\mathrm{P}}\)，使得对于每个 \(\mathfrak{g}\)-模 E 以及 E 中每个权为 \(e\) 的本原元 \(\lambda\)，都有 \(x_{\mathrm{E}}e = \varphi (\lambda)e\)。

（选取 \(\varphi\)，使得当 E 是基本模 \(\mathrm{E}(\varpi_{\alpha})\) 时该关系成立。利用将这类模嵌入若干 \(\mathrm{E}(\lambda)\) 的张量积，推出 \(\lambda \in \mathrm{P}_{++}\)、\(\mathrm{E}(\varpi_{\alpha})\) 的情形。再由此推广到一般情形。）

\(c_{3}\) ) 按 \(\varphi\) ) 选取 \(c_{2}\)。令 E 为最高权为 \(\lambda\) 的单 g-模，令 \(\mu\) 为 E 的一个权；则 \(\lambda - \mu \in Q\)。证明 \(x_{E}\) 在 \(E^{\mu}\) 上的限制是比值为 \(\varphi(\lambda)\psi(\mu - \lambda)\) 的伸缩。（方法同 \(c_{1}\) )。）

\(c_{4}\) ) 若 \(\lambda, \mu \in P_{++}\)，且 \(\alpha \in B\) 不与 \(\lambda\) 或 \(\mu\) 中任一个正交，则

\[
\varphi (\lambda + \mu - \alpha) = \varphi (\lambda + \mu) \psi (- \alpha),
\]

所以 \(\varphi (\alpha) = \psi (\alpha)\)。（使用 \(c_{2}\) )、\(c_{3}\) ) 以及 \(\operatorname {E}(\lambda +\mu -\alpha)\) 到 \(\operatorname {E}(\lambda)\otimes \operatorname {E}(\mu)\) 中的嵌入，参见~习题17。）

\(c_{5})\) 由 \(c_{4})\) 推出 \(\varphi|\mathbf{Q}=\psi\)，并使用 \(c_{3})\) 推出 \(x=t(\varphi)\)。\(d)\) 通过 \(\mathrm{T}_{\mathrm{Q}}\) 将 \(\mathrm{Aut}_{0}(\mathfrak{g})\) 与 \(f\) 的一个子群相识别。由 \(a)\)，有

\[
\operatorname{Aut} _ {e} (\mathfrak {g}) \subset v (\mathrm{G}) \subset \operatorname{Aut} _ {0} (\mathfrak {g})
\]

且由 \(c\) )，\(v(\mathrm{G}) \cap \mathrm{T}_{\mathrm{Q}} = \mathrm{Im}(\mathrm{T}_{\mathrm{P}})\)。推出（参见~§5第3节）\(\mathrm{Aut}_e(\mathfrak{g}) \cap \mathrm{T}_{\mathrm{Q}} = \mathrm{Im}(\mathrm{T}_{\mathrm{P}})\)，且 \(f(\mathrm{G}) = \mathrm{Aut}_e(\mathfrak{g})\)。典则映射

\[
\iota : \mathrm{T} _ {\mathrm{Q}} / \operatorname{Im} (\mathrm{T} _ {\mathrm{P}}) \to \operatorname{Aut} _ {0} (\mathfrak {g}) / \operatorname{Aut} _ {e} (\mathfrak {g})
\]

因此是同构。

\begin{enumerate}
\def\labelenumi{\alph{enumi})}
\setcounter{enumi}{4}
\tightlist
\item
  \(v: \mathrm{G} \to \mathrm{Aut}_e(\mathfrak{g})\) 的核等于 \(\mathrm{T}_{\mathrm{P}} \to \mathrm{T}_{\mathrm{Q}}\) 的核；它同构于
\end{enumerate}

\[
\operatorname{Hom} (\mathrm{P} / \mathrm{Q}, k ^ {*});
\]

这是包含于 G 中心的有限阿贝尔群，其阶整除 \((P:Q)\)；若 k 是代数闭的，则它同构于 P/Q 的对偶（《代数》第VII章§4第8节）。

\(f)\) 令 \(\alpha \in \mathrm{R}\)，\(X_{\alpha} \in \mathfrak{g}^{\alpha}\) 和 \(X_{-\alpha} \in \mathfrak{g}^{-\alpha}\) 满足 \([X_{\alpha}, X_{-\alpha}] = -H_{\alpha}\)，并令 \(\rho_{\alpha}\) 为 \(\mathfrak{sl}(2, k)\) 上相应的 \(\mathfrak{g}\) 表示。若 E 是 \(\mathfrak{g}\)-模，则由§1第4节得到 \(\mathbf{SL}(2, k)\) 在 E 上的表示，因而（由习题25 b)、c)）得到同态

\[
\varphi_ {\alpha}: \mathbf {SL} (2, k) \to \mathrm{G}.
\]

  证明 \(\operatorname{Im}(\varphi_{\alpha})\) 包含 \(T_{P}\) 中形如 \(\lambda \mapsto t^{\lambda(H_{\alpha})}, t \in k^{*}\) 的元素。推出各个 \(\operatorname{Im}(\varphi_{\alpha}), \alpha \in B\) 生成 G（并因而生成 \(T_{P}\)）。特别地，G 由 \(e^{n}\) 生成，其中 \(n \in g^{\alpha}, \alpha \in B \cup -B\)。

\(g\) ) 若 \(\mathbf{G}'\) 是 \(\mathbf{G}\) 的一个子群，满足 \(v(\mathrm{G}') = \mathrm{Aut}_e(\mathfrak{g})\)，则 \(\mathrm{G}' = \mathrm{G}\)（使用 \(f\)）。\(h\) ) 令 \(\mathrm{E}\) 为忠实的 \(\mathfrak{g}\)-模，令 \(\Gamma\) 为 \(\mathrm{P}\) 的权生成的 \(\mathrm{E}\) 的子群。则 \(\mathrm{P} \supset \Gamma \supset \mathrm{Q}\)，参见~习题5。证明典则同态 \(\mathrm{G} \to \mathbf{GL}(\mathrm{E})\) 的核等于 \(T_{P}\) 的子群，该子群由限制在 \(\varphi\) 上为平凡的元素 \(\varGamma\) 组成。特别地，若 \(\varGamma = P\)，同态 \(\mathrm{G} \to \mathbf{GL}(\mathrm{E})\) 是单射。若 \(\varGamma = Q\)，该同态分解为 \(\mathrm{G} \stackrel{v}{\longrightarrow} \operatorname{Aut}_{e}(\mathfrak{g}) \longrightarrow \mathbf{GL}(\mathrm{E})\)，且同态 \(\operatorname{Aut}_{e}(\mathfrak{g}) \to \mathbf{GL}(\mathrm{E})\) 是单射。

\begin{enumerate}
\def\labelenumi{\arabic{enumi})}
\setcounter{enumi}{26}
\tightlist
\item
  令 \(\Omega = \mathrm{P} / \mathrm{Q}\)。若 \(\omega \in \Omega\)，且 \(\mathrm{E}\) 是 \(\mathfrak{g}\)-模，则以 \(\mathrm{E}_{\omega}\) 表示 \(\mathrm{E}^{\lambda}\) 时的 \(\lambda \in \omega\) 的直和。有 \(\mathrm{E} = \bigoplus_{\omega \in \Omega} \mathrm{E}_{\omega}\)。
\end{enumerate}

\begin{enumerate}
\def\labelenumi{\alph{enumi})}
\tightlist
\item
  证明 \(E_{\omega}\) 是 E 的 g-子模。有 \((\mathrm{E}^{*})_{\omega} = (\mathrm{E}_{-\omega})^{*}\) 且
\end{enumerate}

\[
(\mathrm{E} \otimes \mathrm{F}) _ {\omega} = \bigoplus_ {\alpha + \beta = \omega} \mathrm{E} _ {\alpha} \otimes \mathrm{F} _ {\beta}
\]

若 F 是另一个 g-模。

\begin{enumerate}
\def\labelenumi{\alph{enumi})}
\setcounter{enumi}{1}
\item
  令 \(\chi\in\mathrm{Hom}(\varOmega,k^{*})=\mathrm{Ker}(\mathrm{T}_{\mathrm{P}}\to\mathrm{T}_{\mathrm{Q}})\)。将 \(\chi\) 与 \(f:\mathrm{G}\to\mathrm{Aut}_{e}(\mathfrak{g})\) 的核中的一个元素相识别，参见~习题26 e)。证明 \(\chi\) 在 \(E_{\omega}\) 上的作用是比值为 \(\chi(\omega)\) 的伸缩。
\item
  当 \(\mathrm{E}_{\omega}\) 时，\(\mathfrak{g} = \mathfrak{sl}(2,k)\) 是什么？
\end{enumerate}

